\documentclass[../../main.tex]{subfiles}
\begin{document}


{\bf[解]} 
まず，$O$を原点とし，半直線$OX$が$x$軸正方向になるように$xy$座標を定める．このとき$OY$が第$1$象限にあるようにする．さらに動点が時刻$0$に各点を出発したとする．
すると時刻$t$での動点の位置は以下のように表せる．
     \begin{align}
     &P(t+1,0) \\
     &Q(\dfrac{3}{2}t,\dfrac{\sqrt{3}}{2}t) \\
     & R(t+\dfrac{1}{2},\sqrt{3}t+\dfrac{\sqrt{3}}{2})
     \end{align}
これを用いて問に答える．

\paragraph{(1)}
$\vec{PQ}\parallel\vec{PR}$となればよい．故に$k\in\mathbb{R}$として
\begin{align}
  \vec{PQ}=&k\vec{PR} \\
  \begin{pmatrix}
   \dfrac{1}{2}t-1 \\
   \dfrac{\sqrt{3}}{2}t
  \end{pmatrix}
  =&k
  \begin{pmatrix}
   -\dfrac{1}{2} \\
   \sqrt{3}t+\dfrac{\sqrt{3}}{2}
  \end{pmatrix} 
\end{align}
このような$k$の存在条件を求めればよい．ゆえに$k$を消去して
\begin{align}
 \dfrac{\sqrt{3}}{2}t=&(2-t)(\sqrt{3}t+\dfrac{\sqrt{3}}{2}) \\
 t=&\dfrac{1+\sqrt{5}}{2}
\end{align}
よってもとめる値は$t=\dfrac{1+\sqrt{5}}{2}\cdots$(答) である．


\paragraph{(2)}
$\triangle AOB$の面積は$\dfrac{\sqrt{3}}{4}\,\mathrm{cm^2}$であるから，時刻$t$での
$\triangle PQR$の面積がこれに等しい時，
\begin{align}
          \dfrac{\sqrt{3}}{4}=&\dfrac{1}{2}\left|\left(\dfrac{1}{2}t-1\right)\sqrt{3}\left(t+\dfrac{1}{2}\right)%
          +\dfrac{1}{2}\dfrac{\sqrt{3}}{2}t\right| \\
          = &\dfrac{\sqrt{3}}{4}\left|t^2-t-1\right| 
\end{align}
これを解いて
\begin{align}
 &\left|t^2-t-1\right|=1 \\
 &\therefore  t=1,2 \cdots\text{(答)}
\end{align}
となる．
     
\newpage
\end{document}